% !TeX root = ../colloquium-2.tex

% Формулировка достаточного условия экстремума.
% Определение Определение выпуклой функции.
% Формулировка критерия выпуклости.

\phantomsection\label{sec_Definition18}
\begin{definition}[
    Формулировка достаточного условия экстремума.
    Определение выпуклой функции.
    Формулировка критерия выпуклости.
]
    \subsection*{Формулировка достаточного условия экстремума.}
    \begin{note}[purple]
        Пусть функция $f$ определена в некоторой окрестности точки $a$ и дифференцируема $n$ 
        раз в точке $a$, причём 
        $$ f'(a) = ... = f^{(n - 1)}(a) = 0, $$
        а $f^{(n)}(a) \neq 0$. Тогда если $n = 2k$ и $f^{(n)}(a) < 0$, то $a$ - точка максимума,
        если $n = 2k$ и $f^{(n)}(a) > 0$, то $a$ - точка минимума, а если $n = 2k - 1$, то
        экстремума нет.
    \end{note}

    \subsection*{Определение выпуклой функции.}
    \begin{note}[purple]
        \subsubsection*{Определение выпуклости}
        Функция $f$, определённая на интервале $(a, \; b)$, называется выпуклой 
        (\textit{выпуклой вниз}) на этом интервале, если для любых точек $x_1, x_2 \in (a, \; b)$
        и любого числа $\lambda \in [0, \; 1]$ выполнено неравенство
        $$ f(\lambda x_1 + (1 - \lambda)x_2) \leq \lambda f(x_1) + (1 - \lambda) f(x_2) $$

        Если при тех же условиях выполнено противоположное неравенство, то функцию $f$ называют
        вогнутой (\textit{выпуклой вверх}) на интервале $(a, \; b)$/

        \subsubsection*{Геометрическое определение выпуклости}
        Всякая хорда, соединяющая две точки граифика функции, лежит не ниже графика функции,
        а определение вогнутости означает, что любая хорда лежит не выше графика функции.
    \end{note}

    \subsection*{Формулировка критерия выпуклости.}
    \begin{note}[purple]
        Функция $f$ выпукла на интервале $(a, \; b)$ тогда и только тогда, когда для любых
        таких чисел $x_1, \; x_2, \; y \in (a, \; b)$, что $x_1 < y < x_2$, выполнено 
        неравенство 
        $$ 
        \frac{f(y) - f(x_1)}{y - x_1} \leq 
        \frac{f(x_2) - f(y)}{x_2 - y}
        $$
        Геометрический смысл заключается в том, что тангенсы углов наклонов секущих не убывают
        с увеличением абсцисс точек пересечения с графиком функции $f$. 
    \end{note}
\end{definition}